\chapter{Navegación aérea}\label{ch:14}
\objetivos{Separar posición, navegación, guiado y control; utilizar referencias geográficas y angulares coherentes; resolver el triángulo de velocidades con viento; explicar posicionamiento inercial, radioeléctrico y satelital; interpretar error, integridad y límites de las ayudas.}

\section{Saber dónde se está y decidir hacia dónde ir}
Navegar consiste en estimar posición y movimiento respecto a un sistema de referencia y relacionarlos con una ruta. El guiado convierte una ruta deseada en consignas de trayectoria o actitud. El control actúa sobre la aeronave para seguir esas consignas. Una medida precisa de posición no garantiza por sí sola que el vehículo siga una trayectoria segura: también hacen falta una ruta bien definida, estimación de estado, lógica de guiado y capacidad de mando.

La posición responde a ``dónde'', la velocidad a ``cómo cambia la posición'' y el tiempo permite sincronizar ambas. Una solución navegacional incluye incertidumbre y antigüedad. Dos posiciones obtenidas con igual precisión espacial pueden ser muy diferentes para guiar una aeronave rápida si una llega con retraso. En una introducción universitaria conviene analizar observaciones, modelo y decisión como etapas relacionadas pero distintas. FAA presenta métodos básicos de navegación y su relación con planificación y viento \cite{PHAK}.

\section{Coordenadas y aproximaciones geográficas}
La latitud $\varphi$ expresa posición norte o sur y la longitud $\lambda$ posición este u oeste respecto al meridiano de referencia. La Tierra se aproxima mejor mediante un elipsoide que mediante una esfera, pero una esfera permite derivar relaciones introductorias. La altura debe indicar referencia: superficie elipsoidal, nivel medio del mar o un ajuste barométrico no representan exactamente la misma magnitud.

Para desplazamientos pequeños sobre una esfera de radio $R$, las componentes locales son
\begin{equation}\label{eq:14-local}
\Delta N\simeq R\Delta\varphi,\qquad
\Delta E\simeq R\cos\varphi\,\Delta\lambda,
\end{equation}
con ángulos en radianes. La distancia por grado de longitud disminuye al aumentar latitud. Olvidar el factor coseno equivale a tratar todos los paralelos como si fueran el ecuador. Cerca de los polos y para tramos largos se requiere una representación geodésica apropiada.

En una esfera, la distancia más corta entre puntos se encuentra sobre un arco de círculo máximo. Si $\sigma$ es el ángulo central,
\begin{equation}
\sin^2\frac\sigma2=\sin^2\frac{\Delta\varphi}2+
\cos\varphi_1\cos\varphi_2\sin^2\frac{\Delta\lambda}2,
\qquad d=R\sigma.
\end{equation}
La relación evita confundir diferencias angulares con distancias. Una loxodromia corta los meridianos con ángulo constante y resulta cómoda para mantener rumbo sobre ciertas proyecciones, pero salvo casos particulares no minimiza distancia esférica. La ruta operativa puede diferir de ambas por espacios disponibles, meteorología, prestaciones y restricciones publicadas.

La milla náutica internacional equivale exactamente a $1852\ \mathrm m$ y el nudo a una milla náutica por hora. Un vuelo a $120$ nudos no recorre $120$ kilómetros por hora. Las conversiones deben realizarse antes de combinar datos: $1\ \mathrm{kt}=1852/3600\ \mathrm{m/s}$. En los ejemplos usaremos unidades SI para que la suma de velocidades sea transparente.

\section{Rumbo, derrota y referencias de norte}
El rumbo es la dirección hacia la que apunta el eje longitudinal; la derrota es la dirección del movimiento sobre el terreno. Con viento cruzado difieren. El curso deseado describe una dirección planificada, mientras que una derrota observada describe lo que está ocurriendo. La deriva es la diferencia angular correspondiente según la convención elegida. En un dibujo debe indicarse qué ángulo se mide desde qué dirección y en qué sentido.

Los ángulos verdaderos se refieren al norte geográfico; los magnéticos, al norte magnético. Si la declinación $D$ se considera positiva hacia el este, se cumple $\psi_M=\psi_T-D$, tomando después el ángulo módulo $360^\circ$. Si la desviación de brújula $\delta_c$ se define con la misma convención, $\psi_C=\psi_M-\delta_c$. Las palabras ``sumar este'' sin una definición pueden inducir un signo incorrecto. La declinación depende del lugar y de la fecha; la desviación depende además de la instalación y del rumbo.

Un cálculo de viento no debe mezclar una dirección verdadera con otra magnética. Tampoco debe utilizar una designación de pista redondeada cuando necesita orientación precisa. Antes de resolver el problema, se convierten todas las direcciones a una referencia común. Para operaciones reales se consulta información vigente; este capítulo deriva geometría y no proporciona rumbos publicados de rutas concretas.

\section{El triángulo de velocidades}
La relación fundamental entre velocidad respecto al suelo, velocidad respecto al aire y viento es vectorial:
\begin{equation}\label{eq:14-vector}
\boldsymbol V_g=\boldsymbol V_a+\boldsymbol V_w.
\end{equation}
La velocidad del viento apunta hacia donde se desplaza el aire. La dirección meteorológica indica de dónde procede: un viento del norte tiene componente norte negativa. Con componentes este y norte y un rumbo verdadero $\psi$ medido desde el norte en sentido horario,
\begin{equation}
V_{a,E}=V_a\sin\psi,\qquad V_{a,N}=V_a\cos\psi.
\end{equation}
Estas ecuaciones evitan dibujar el triángulo con una orientación implícita.
\figura{wind}{Suma de velocidad aérea y viento para seguir una derrota deseada.}

Si se desea una derrota de dirección unitaria $\boldsymbol t$ y normal lateral $\boldsymbol n$, la condición de no derivar lateralmente es
\begin{equation}
\boldsymbol V_a\cdot\boldsymbol n=-\boldsymbol V_w\cdot\boldsymbol n.
\end{equation}
Existe una solución geométrica hacia delante sólo si la componente transversal de viento no supera la rapidez aérea y la componente longitudinal resultante sobre el suelo es positiva. Un valor de arcseno real no garantiza por sí solo avance en el sentido deseado. En el límite de viento transversal igual a la rapidez aérea, toda la velocidad aérea se consume compensando lateralmente y puede no quedar progreso útil.

La magnitud de velocidad sobre el suelo determina el tiempo: $t=d/V_g$ para un tramo recto con viento y velocidad constantes. El consumo depende del tiempo y del régimen propulsivo, por lo que una diferencia de viento puede afectar autonomía aunque la trayectoria geométrica sea idéntica. Para viento variable se integran tramos o se resuelve una evolución; utilizar la media de direcciones sin componentes puede dar un viento medio incorrecto.

\section{Estima y navegación observada}
La navegación a estima propaga una posición conocida usando rumbo, velocidad y tiempo, junto con una estimación de viento. Es sencilla y permite una comprobación independiente de otros sistemas, pero acumula errores. Una desviación angular pequeña mantenida durante distancia $s$ produce un error transversal aproximadamente $s\Delta\psi$, con $\Delta\psi$ en radianes. Así, un grado durante cien kilómetros genera cerca de $1{,}75\ \mathrm{km}$ de separación lateral.

La navegación visual compara el terreno con referencias cartográficas. La calidad depende de visibilidad, identificación inequívoca y distribución de puntos. Confundir una referencia puede producir un error sistemático mayor que el de una medida ruidosa. Varias observaciones independientes ayudan a detectar incoherencias; varias observaciones del mismo punto mal identificado pueden reforzar el mismo error.

Una línea de posición restringe la ubicación a un conjunto de puntos: una distancia a una estación define una circunferencia ideal en el plano y una dirección define una recta o radial. Dos líneas apropiadamente orientadas pueden proporcionar una intersección. Si se cortan con ángulo muy pequeño, el error de posición aumenta: la geometría de las observaciones importa tanto como la precisión de cada sensor.

\section{Radioayudas y guiado de aproximación}
Un VOR proporciona información angular respecto a una estación. La radial identifica la dirección desde la estación hacia el avión; no equivale sin más al rumbo para volar hacia ella. Un NDB, con un receptor de radiogoniometría, permite estimar dirección relativa, pero la señal y su interpretación pueden estar afectadas por propagación y condiciones de recepción. Se deben distinguir referencia angular, sentido hacia o desde la estación y trayectoria deseada.

El DME mide una distancia oblicua por tiempo de propagación de señales, no una distancia horizontal exacta. Si la diferencia de altura es $h$ y la separación horizontal es $d$, el modelo geométrico elemental da $R_{\mathrm{DME}}^2=d^2+h^2$. Sobre la estación, la distancia horizontal puede ser casi nula y la lectura seguir siendo aproximadamente la altura relativa. Lejos de ella y con $h\ll d$, ambas distancias se aproximan.

El ILS proporciona señales de guiado lateral y vertical mediante localizador y senda. Una indicación centrada describe relación con esas señales; su uso depende de procedimiento, disponibilidad y prestaciones del conjunto. Una ayuda de guiado no aporta automáticamente todas las funciones de posicionamiento global. Las cartas y publicaciones especifican referencias, limitaciones y datos operativos de cada instalación.

\section{Posicionamiento satelital e inercial}
Los sistemas globales por satélite estiman posición mediante observaciones asociadas a propagación y tiempo. En un modelo simplificado, la pseudodistancia al satélite $i$ es
\begin{equation}\label{eq:14-gnss}
\rho_i=\|\boldsymbol r-\boldsymbol r_i\|+c\,\delta t+\varepsilon_i,
\end{equation}
donde $\boldsymbol r$ es posición del receptor, $\boldsymbol r_i$ la del satélite, $\delta t$ el error de reloj del receptor y $\varepsilon_i$ reúne errores no corregidos. El símbolo $c$ representa aquí velocidad de la luz, no cuerda aerodinámica. Tres coordenadas y un desfase de reloj introducen cuatro incógnitas; se requieren al menos cuatro observaciones independientes con geometría adecuada para la solución básica tridimensional.

Más satélites pueden aportar redundancia, pero no eliminan todos los errores. Intervienen órbitas y relojes, propagación atmosférica, multitrayectoria y perturbaciones de señal. Los sistemas de aumentación aportan correcciones y elementos de integridad. Precisión es proximidad al valor verdadero; integridad incluye capacidad de advertir cuando la solución no debe utilizarse bajo un criterio establecido. Una solución aparentemente suave puede ser incorrecta si existe un sesgo común.

La navegación inercial usa giróscopos y acelerómetros para propagar actitud, velocidad y posición. Los acelerómetros miden fuerza específica; para obtener aceleración respecto al marco elegido hay que considerar gravedad y transformaciones de referencia. Un sesgo constante de aceleración no corregido produce, en un modelo simple, un error de velocidad proporcional al tiempo y un error de posición aproximadamente $bt^2/2$.

La integración de sistemas puede combinar continuidad inercial con referencias externas. La fusión debe ponderar incertidumbres y detectar incoherencias; no consiste en promediar números sin modelo. Una observación satelital errónea puede contaminar la estimación si se acepta como cierta. La redundancia útil depende de independencia, geometría y capacidad de diagnóstico.

\section{Problema resuelto: compensación de viento cruzado}
\paragraph{1. Enunciado.} Calcule el rumbo para recorrer un tramo exactamente hacia el este y su tiempo de vuelo con viento uniforme del norte.
\paragraph{2. Datos.} Derrota verdadera deseada $090^\circ$, rapidez aérea $55{,}0\ \mathrm{m/s}$, viento del norte de $20{,}0\ \mathrm{m/s}$ y distancia sobre el terreno $100{,}0\ \mathrm{km}$.
\paragraph{3. Incógnitas.} Rumbo verdadero $\psi$, velocidad sobre el suelo y tiempo.
\paragraph{4. Esquema verbal.} El viento desplaza al avión hacia el sur. El morro debe apuntar al norte de la derrota para que la componente norte de velocidad aérea lo compense.
\paragraph{5. Hipótesis.} Tierra localmente plana, altura constante, viento uniforme y todas las direcciones referidas al norte verdadero. No se calcula consumo.
\paragraph{6. Modelo físico.} Se suman vectores de velocidad y se anula la componente norte de la velocidad sobre el suelo.
\paragraph{7. Ecuaciones simbólicas.} $V_{g,N}=55\cos\psi-20=0$, $V_{g,E}=55\sin\psi$ y $t=d/V_{g,E}$.
\paragraph{8. Despeje.} $\psi=\arccos(20/55)$ para la solución con avance hacia el este; $V_g=\sqrt{55^2-20^2}$.
\paragraph{9. Sustitución.} $\psi=\arccos(0{,}363636)$ y $t=100000/\sqrt{3025-400}$.
\paragraph{10. Resultado con unidades.} $\psi=68{,}676^\circ$, $V_g=51{,}235\ \mathrm{m/s}$ y $t=1951{,}80\ \mathrm s=32{,}530\ \mathrm{min}$.
\paragraph{11. Interpretación.} La corrección es $21{,}324^\circ$ hacia el norte. La velocidad sobre el suelo es menor que la aérea porque parte de ésta compensa el viento lateral.
\paragraph{12. Comprobación independiente.} Las componentes aéreas son $(51{,}235,20{,}000)\ \mathrm{m/s}$ en orden este--norte. Al sumar viento $(0,-20{,}000)$ queda $(51{,}235,0)$; su módulo aéreo original recupera $55\ \mathrm{m/s}$ por Pitágoras.
\paragraph{13. Limitaciones.} No se aplican declinación ni desviación de brújula porque no se han suministrado. Variaciones de viento o curvatura apreciable de ruta exigirían actualizar rumbo y tiempo.

\section{Ejercicios propuestos y guía de respuestas}
\begin{enumerate}
\item Con rumbo verdadero $100^\circ$ y declinación este de $4^\circ$, calcule rumbo magnético. Respuesta: $96^\circ$, con la convención del texto.
\item Convierta $90\ \mathrm{kt}$ a SI. Respuesta: $46{,}30\ \mathrm{m/s}$.
\item Un DME indica $5\ \mathrm{km}$ y existe diferencia vertical de $3\ \mathrm{km}$. Obtenga separación horizontal ideal. Respuesta: $4\ \mathrm{km}$.
\item Un sesgo de aceleración de $0{,}002\ \mathrm{m/s^2}$ actúa sin corregir durante $100\ \mathrm s$. Estime error de posición unidimensional. Respuesta: $10\ \mathrm m$, con estado inicial exacto.
\item ¿Por qué cuatro observaciones satelitales pueden dar una solución pobre? Guía: analizar geometría, errores comunes y falta de redundancia; cuatro ecuaciones no garantizan una solución robusta.
\end{enumerate}
\section{Resumen y autoevaluación}
\resumen{Navegación estima estado y referencia; guiado define cómo seguir una trayectoria y control ejecuta la respuesta. El viento exige sumar velocidades como vectores. Los métodos observados, inerciales y satelitales aportan información y errores distintos. Referencias angulares, unidades, tiempo e integridad deben acompañar a cualquier posición calculada.}
¿Por qué la derrota difiere del rumbo con viento cruzado? ¿Qué mide una pseudodistancia que no sea una distancia geométrica pura? ¿Qué diferencia precisión e integridad? ¿Por qué una ayuda angular y una distancia oblicua requieren un modelo geométrico para fijar posición?
